% Generated draft: requires mathematical and source review.
\documentclass[11pt]{article}
\usepackage[margin=25mm]{geometry}
\usepackage{fontspec}
\setmainfont{Latin Modern Roman}
\usepackage{amsmath,amssymb,mathtools,unicode-math}
\setmathfont{Latin Modern Math}
\usepackage{enumitem,needspace,xurl,hyperref}
\setlength{\parindent}{0pt}
\setlength{\parskip}{5pt}
\setlength{\emergencystretch}{3em}
\newcommand{\N}{\mathbb{N}}
\newcommand{\Z}{\mathbb{Z}}
\newcommand{\Q}{\mathbb{Q}}
\newcommand{\R}{\mathbb{R}}
\newcommand{\C}{\mathbb{C}}
\newcommand{\F}{\mathbb{F}}
\newcommand{\A}{\mathbb{A}}
\newcommand{\PP}{\mathbb{P}}
\newcommand{\E}{\mathbb{E}}
\newcommand{\eps}{\varepsilon}
\newcommand{\dd}{\,\mathrm{d}}
\begin{document}
\section*{UnsolvedMath ID: 30003042}
\textbf{Source problem code:} catalogue entry
\subsection{Definitions and mathematical statement}
Let \(n\) be an integer with \(n\geq 8\). The supplied record uses the notation \(K(X_{F(S)})\) and calls \(e_n\) a special Schubert class in this group. It does not specify the field \(F\), the space \(S\), the variety \(X\), their dependence on \(n\), or the construction of \(e_n\).

For clarity about the filtration convention, if the intended \(K\)-theory is a Grothendieck \(\lambda\)-ring \(R\) of vector bundles equipped with a rank augmentation \(\operatorname{rk}\colon R\to\Z\), put \(I=\ker(\operatorname{rk})\). Define \(\gamma^i(x)\) by
\[
\sum_{i\geq 0}\gamma^i(x)t^i=\lambda_{t/(1-t)}(x),\qquad x\in I.
\]
Under this convention, for \(r\geq 1\), the subgroup \(\gamma^rR\) is generated additively by products
\[
\gamma^{i_1}(x_1)\cdots\gamma^{i_m}(x_m),\qquad
m\geq 1,\quad x_j\in I,\quad i_j\geq 1,\quad \sum_{j=1}^m i_j\geq r.
\]
The record does not explicitly identify its \(K\)-theory or augmentation with this convention.

\textbf{Supplied conjectural assertion.} With the intended geometric setup and special Schubert class from the cited source, the assertion is that, for every integer \(n\geq 8\),
\[
e_n\in K(X_{F(S)})\qquad\text{and}\qquad
e_n\notin\gamma^{\,n+1-\lfloor(n+1)/2\rfloor}K(X_{F(S)}).
\]
The omitted geometric setup and the normalization of \(e_n\) must be recovered from the cited source before this can be treated as a complete standalone mathematical statement. No unrestricted quantification over fields or varieties is asserted here.
\subsection{Short English statement}
The supplied conjecture says that, for every integer \(n\geq 8\), a designated special Schubert class lies outside the \(n+1-\lfloor(n+1)/2\rfloor\) term of the gamma filtration. The underlying geometric objects and the class are not defined in the record, so the assertion remains dependent on its source context.
\subsection{Sources}
\begin{itemize}
\item Local UnsolvedMath registry, supplied record with permanent ID \(30003042\), title \emph{Schubert Classes in the Gamma Filtration}, exact source problem code \texttt{catalogue entry}; Oberwolfach Reports Open Problems collection, proposed year \(2016\). No proposer is identified in the supplied record.
\item \emph{Algebraic Cobordism and Projective Homogeneous Varieties} (2016), Oberwolfach Reports, pp.~215--265, DOI \texttt{10.4171/owr/2016/5}. Supplied report reference and literature-source link labeled ``original OWR report'': \url{https://doi.org/10.4171/owr/2016/5}. The supplied background describes these reports as documenting workshops and their open-problem discussions.
\item Supplied literature source labeled ``Karpenko on 16-dimensional forms'': \url{https://sites.ualberta.ca/~karpenko/publ/a16i3.pdf}. The record supplies no fuller bibliographic citation for this item.
\item Literature assessment dated exactly \texttt{2026-08-21}: the supplied curation labels the problem open and reports that the special Schubert-class gamma-filtration assertion remains a conjecture in Karpenko's primary treatment, with no verified later proof of this exact claim found. This is supplied dated curation derived from the OWR curation, not an independently verified literature assessment; its status and cited claims require independent verification.
\item The registry labels the statement \texttt{corrected\_verified}, with the qualification that only safe notation, typography, and TeX-delimiter cleanup was applied and all hypotheses, quantifiers, and subquestions were preserved. This supplied label does not establish mathematical completeness or correctness.
\end{itemize}
\end{document}
